\documentclass[10pt,a4paper]{amsart}
\usepackage[margin=19mm]{geometry}
\usepackage{amsmath,amssymb,mathtools}
\usepackage[scr=rsfs,cal=euler]{mathalfa}
\usepackage{xfrac}
\usepackage[unicode,hidelinks,bookmarks=false]{hyperref}
\newcommand{\ie}{i.e.\ }
\newcommand{\cf}{cf.\ }
\newcommand{\resp}{respectively\ }
\newcommand{\Subsection}{\subsection}
\providecommand{\nobreakdash}{\nobreak\hbox{-}}
\setlength{\emergencystretch}{2em}
\multlinegap0pt
\usepackage{amssymb}
\usepackage{mathtools}
\usepackage{bm}
\usepackage{xcolor}
\newtheorem{lem}{Lemma}
\newcommand{\lan}{\langle }
\newcommand{\mc}{\mathcal}
\newcommand{\mcr}{\mathscr}
\newcommand{\mkngam}{M^n_k(\Gamma)}
\newcommand{\q}{\quad}
\newcommand{\ran}{\rangle}
\title{Fourier coefficients and cuspidality for modular forms: another approach and beyond}
\author{Soumya Das}
\date{Source: arXiv:2407.15222v3 (CC BY 4.0)}
\begin{document}
\maketitle

\noindent\emph{Source and licence.} Fourier coefficients and cuspidality for modular forms: another approach and beyond. Version arXiv:2407.15222v3 (CC BY 4.0), distributed by arXiv under the Creative Commons Attribution 4.0 International licence, http://creativecommons.org/licenses/by/4.0/, as recorded in the arXiv licence metadata for this exact version and shown on the abstract page https://arxiv.org/abs/2407.15222v3. Full licence notice: \texttt{licenses/arxiv-2407.15222-CC-BY-4.0.txt}.\par\smallskip

\noindent\textit{Editorial scope.} Counted unit: the article's lem polu (source0.tex lines 752--760). The article's complete original proof of it is delivered with it (source0.tex lines 762--772, one \texttt{proof} environment of 11 lines). No mathematics is written, repaired or re-derived: the statement and the proof are the article's own lines, in the article's own notation, and no retained line is substituted or re-flowed.  Retained uncounted as the source-local context the proof needs: lines 699--704; lines 711--717; lines 742--746.  Nothing was omitted from any retained span, so the package carries no omission record.  No retained line contains a citation, so the wrapper carries no bibliography and no bibliography entry.  No retained line contains a figure environment, an image-inclusion command or drawing code, so no image is delivered and the package declares no asset dependency.  Editorial formatting only, in addition: an AMS \texttt{amsart} preamble that restates the article's own declarations and macros used by the retained lines, namely the environments \texttt{lem} on separate counters, together with the standard packages those lines need.  The retained environments print in this document as Lemma 1 in order of appearance, whereas the article prints the same items with the numbering of its own full document.  The complete native source of the article as distributed in the arXiv e-print for this exact version is bundled unchanged as \texttt{sources/162646/source0.tex}.
\begin{align} \label{rf-diffop}
\mc R^*(F,s) := \begin{cases}
\lan E^n(s +\frac{n+1}{2}-k, \cdot), \, \det(Y)^{k}|F|^2 \ran \q &\text{ if } F \text{ is a cusp form}, \\
\lan E^n(s +\frac{n+1}{2}-k, \cdot), \, \mc D_n(\det(Y)^{k}|F|^2)  \ran \q &\text{ if } F \text{ is a not cusp form}. 
\end{cases}
\end{align}

\begin{align} \label{rf-int-rep}
 \mc R^*(F,s) =\begin{cases}
 c_n(s) \Gamma_n(s) \mc R(F,s)
 \q &\text{ if } F \text{ is a cusp form}, \\
  c_n(s) \Gamma_n(s) \phi_n(s) \phi_n(s+n-2k)  \mc R(F,s) \q &\text{ if } F \text{ is a not cusp form}. 
\end{cases}
\end{align}

\begin{lem} \label{pole-eisen}
For the weight $0$ (completed) Eisenstein series $\mathscr E^n(s,Z) $, we have,
\[ \mathscr{P}( \mathscr E^n(s,Z) ) \subseteq \{  j/4, \, 0 \le j \le 2n+2 \} . \]
Moreover its right-most pole  is at $s=\frac{n+1}{2}$. It is simple with residue $\prod_{j=1}^{[n/2]} \xi(2j+1)$.
\end{lem}

\begin{lem} \label{polu}
With the above notation, let $F \in \mkngam$ and $k>n$. Then,
\begin{align}
    \mathscr{P}(\mc R(F,s)) \subseteq \begin{cases}
       \mathscr{A}  \q & F \text{ is a cusp form} \\
       \mathscr{A}  \cup \mathscr{C} \q & F \text{ is not a cusp form}.
    \end{cases}
\end{align}
\end{lem}

\begin{proof}
    When $F$ is a cusp form, the Lemma follows from \eqref{rf-diffop}, \eqref{rf-int-rep} and the knowledge of poles of $\mc R^*(F,s)$, viz. at $t_j=k-j/4, \,   0 \le j \le 2n+2$ as mentioned above. More precisely, if we consider \eqref{rf-diffop}, we see that $\mathscr{P}(\mc R^*(F,s)) = \mathscr{P}( \mathscr E^n(s +\frac{n+1}{2}-k,Z) )$. 
    
    Then from \eqref{rf-int-rep}, it follows that 
    \[ \mathscr{P}(\mc R(F,s)) = \mathscr{P}( \Gamma_n(s)^{-1} \mc R^*(F,s)) =\mathscr{P}(\mc R^*(F,s)) . \]
    To see the second equality above, note that since $k>n$, the zeros of $\Gamma_n(s)^{-1}$ do not cancel any of these poles of $\mc R^*(F,s)$. We already know all the possible poles of $\mc R^*(F,s)$ from \eqref{rf-diffop}, as noted in the above paragraph. Therefore, when $F$ is a cusp form, we have,
    \[  \mathscr{P}(\mc R(F,s)) =  \mathscr{P}( \mathscr E^n(s +\frac{n+1}{2}-k,Z) ) \subseteq \mcr A,\]
if we consult Lemma~\ref{pole-eisen}.

    If $F$ is not a cusp form, we have to a bit more careful, but the same idea works. In this case, we repeat the above process, i.e. look at \eqref{rf-diffop} and \eqref{rf-int-rep}. The extra possible poles come from the polynomial $\phi_n(s+n-2k)$ and these constitute the set $\mcr C$. Note that the zeros of $\Gamma_n(s)^{-1} \phi_n(s)^{-1}$ do not cancel any of these possible poles of $\mc R^*(F,s)$.
\end{proof}

\end{document}
